\documentclass{article}

\usepackage{ctex}
\usepackage{amsmath, amssymb, amsthm, amsfonts}
\newtheorem{definition}{Definition}
\newtheorem*{note}{Note}
\newtheorem*{remark}{Remark}
\newtheorem*{example}{Example}
\newtheorem*{principle}{Principle}
\newtheorem*{corollary}{corollary}
\newtheorem{theorem}{定理}

\usepackage{amsmath,amssymb}
\usepackage{geometry}
\usepackage{enumitem}
\usepackage{xcolor}
\usepackage{booktabs}
\usepackage{ctex}
\usepackage{tikz}
\usepackage{enumitem}
\usepackage{hyperref}

\begin{document}


\section{数学历史}

\href{https://mathshistory.st-andrews.ac.uk/}{一个好的数学网站}

\section{二次型与圆锥曲线分类}

\subsection{圆锥曲线的一般形式}

平面上任意一条圆锥曲线都可以表示为如下二次方程：
\[
Ax^2 + Bxy + Cy^2 + Dx + Ey + F = 0,
\]
其中 \(A,B,C,D,E,F \in \mathbb{R}\)。  
其中的二次项部分
\[
Q(x, y) = Ax^2 + Bxy + Cy^2
\]
定义了一个二次型（Quadratic Form）。


\subsection{二次型的矩阵形式}

我们可以用矩阵的形式将二次型写作：
\[
Q(x, y) =
\begin{bmatrix}
	x & y
\end{bmatrix}
\begin{bmatrix}
	A & \tfrac{B}{2} \\
	\tfrac{B}{2} & C
\end{bmatrix}
\begin{bmatrix}
	x \\ y
\end{bmatrix}
= \mathbf{x}^T M \mathbf{x},
\]
其中
\[
M =
\begin{bmatrix}
	A & \tfrac{B}{2} \\
	\tfrac{B}{2} & C
\end{bmatrix}
\]
是一个对称矩阵，称为该二次型的矩阵。


\subsection{判别式与圆锥曲线分类}

由二次型矩阵 \(M\) 可定义一个判别式：
\[
\Delta = \det(M) = AC - \frac{B^2}{4}.
\]
其符号决定了圆锥曲线的类型：

\[
\begin{array}{c|c|c}
	\text{类型} & \text{条件} & \text{几何意义} \\ \hline
	\text{椭圆（或圆）} & \Delta > 0,\; A,C\text{ 同号} & \text{封闭曲线} \\
	\text{双曲线} & \Delta < 0 & \text{开口曲线} \\
	\text{抛物线} & \Delta = 0 & \text{过渡形曲线}
\end{array}
\]



\subsection{判别式与特征值的关系}

矩阵 \(M\) 的特征多项式为：
\[
\det(M - \lambda I)
= \lambda^2 - (A + C)\lambda + \left(AC - \frac{B^2}{4}\right) = 0.
\]
因此，
\[
\lambda_1 + \lambda_2 = A + C, \qquad
\lambda_1 \lambda_2 = AC - \frac{B^2}{4} = \Delta.
\]
从而得到判别式与特征值之间的代数关系：
\[
\boxed{\Delta = \lambda_1 \lambda_2.}
\]



\subsection{特征值与曲线形状的对应关系}

\begin{table}[ht]
	\centering
	\begin{tabular}{|c|c|c|c|}
		\hline
		特征值符号 & 判别式 & 几何类型 & 方程典型形式 \\
		\hline
		\(\lambda_1, \lambda_2 > 0\)（或都 < 0） & \(\Delta > 0\) & 椭圆（或圆） & \(\dfrac{x'^2}{a^2} + \dfrac{y'^2}{b^2} = 1\) \\
		\hline
		一正一负 & \(\Delta < 0\) & 双曲线 & \(\dfrac{x'^2}{a^2} - \dfrac{y'^2}{b^2} = 1\) \\
		\hline
		一个为零 & \(\Delta = 0\) & 抛物线 & \(y'^2 = 2p x'\) \\
		\hline
	\end{tabular}
\end{table}

几何意义上，特征值的符号决定了曲线的“弯曲方向”：
\begin{itemize}
	\item 两个方向都弯向同一侧（正定）：椭圆；
	\item 一个方向上弯上，一个方向弯下（不定）：双曲线；
	\item 一个方向平坦（半正定）：抛物线。
\end{itemize}


\subsection{通过旋转消去交叉项 \texorpdfstring{\(Bxy\)}{Bxy}}

由于 \(M\) 为实对称矩阵，可被正交矩阵 \(P\) 对角化：
\[
P^T M P = \Lambda =
\begin{bmatrix}
	\lambda_1 & 0 \\ 0 & \lambda_2
\end{bmatrix}.
\]
令
\[
\begin{bmatrix}
	x \\ y
\end{bmatrix}
=
P
\begin{bmatrix}
	x' \\ y'
\end{bmatrix},
\]
则
\[
Q(x, y) = \lambda_1 x'^2 + \lambda_2 y'^2.
\]
这意味着通过坐标旋转可以消去 \(Bxy\) 项。  
旋转角满足：
\[
\tan 2\theta = \frac{B}{A - C}.
\]


\subsection{平移消去一次项 \texorpdfstring{\(Dx + Ey\)}{Dx + Ey}}

旋转之后，方程形如：
\[
A'x'^2 + C'y'^2 + D'x' + E'y' + F = 0.
\]
设平移：
\[
x' = X + h,\quad y' = Y + k.
\]
展开并消去一次项要求：
\[
\begin{cases}
	2A'h + D' = 0,\\
	2C'k + E' = 0,
\end{cases}
\Rightarrow
h = -\dfrac{D'}{2A'}, \quad k = -\dfrac{E'}{2C'}.
\]
代入后得：
\[
A'X^2 + C'Y^2 + F' = 0.
\]


\subsection{标准方程形式}

1. 椭圆
\[
\frac{X^2}{a^2} + \frac{Y^2}{b^2} = 1, \quad
a^2 = -\frac{F'}{A'},\quad b^2 = -\frac{F'}{C'}.
\]

2. 双曲线
\[
\frac{X^2}{a^2} - \frac{Y^2}{b^2} = 1, \quad
a^2 = \frac{F'}{A'},\quad b^2 = -\frac{F'}{C'}.
\]

3. 抛物线
若 \(A' = 0\)，则
\[
Y^2 = 2pX.
\]


\subsection*{九、几何解释：特征值与曲面的形状}

设
\[
Q(x, y) = \lambda_1 x'^2 + \lambda_2 y'^2.
\]
将其看作一个“高度函数”，不同符号组合对应不同几何曲面：

\begin{table}[ht]
	\centering
	\begin{tabular}{|c|c|c|}
		\hline
		特征值结构 & 二次型曲面形状 & 等高线（圆锥曲线）  \\
		\hline
		 \(++\) & 椭球面（碗状） & 椭圆 \\
		\hline
		\(+-\) & 马鞍面 & 双曲线 \\
		\hline
		\(+0\) & 抛物面 & 抛物线 \\
		\hline
	\end{tabular}
\end{table}

换言之，\textbf{圆锥曲线是二次型曲面的等高线，而特征值决定其主曲率方向与类型}。

\subsection*{十、总结}

综上，圆锥曲线的分类可由二次型矩阵的代数性质决定：
\[
\boxed{
	\begin{aligned}
		\text{迹：}&\quad \operatorname{tr}(M) = A + C = \lambda_1 + \lambda_2, \\
		\text{行列式：}&\quad \det(M) = AC - \tfrac{B^2}{4} = \lambda_1 \lambda_2, \\
		\text{符号结构：}&\quad \text{sign}(\lambda_1),\text{sign}(\lambda_2)
		\;\Rightarrow\;
		\text{椭圆 / 双曲 / 抛物}.
	\end{aligned}
}
\]
因此，从二次型的角度看，圆锥曲线的几何分类与线性代数中实对称矩阵的谱结构完全一致。


	\section*{圆锥曲线的统一几何定义}
	
	\subsection*{一、核心定义}
	
	在平面内，给定：
	\begin{itemize}
		\item 一个定点 \( F \)（称为\textbf{焦点}）
		\item 一条定直线 \( l \)（称为\textbf{准线}），且 \( F \notin l \)
		\item 一个正常数 \( e \)（称为\textbf{离心率}）
	\end{itemize}
	
	则平面内所有满足以下条件的点 \( P \) 的集合称为\textbf{圆锥曲线}：
	\[
	\frac{\text{点}P\text{到焦点}F\text{的距离}}{\text{点}P\text{到准线}l\text{的距离}} = e
	\]
	用数学符号表示为：
	\[
	\frac{|PF|}{d(P, l)} = e
	\]
	其中 \( d(P, l) \) 表示点 \( P \) 到直线 \( l \) 的垂直距离。
	
	\subsection*{二、分类与离心率}
	
	圆锥曲线的具体形状由离心率 \( e \) 的值唯一确定：
	
	\[
	\text{曲线类型} =
	\begin{cases}
		\text{椭圆（Ellipse）}, & \text{当 } 0 < e < 1 \\
		\text{抛物线（Parabola）}, & \text{当 } e = 1 \\
		\text{双曲线（Hyperbola）}, & \text{当 } e > 1
	\end{cases}
	\]
	
	\begin{center}
		\begin{tikzpicture}[scale=0.8]
			% 坐标轴
			\draw[->] (-1,0) -- (5,0) node[right] {$x$};
			\draw[->] (0,-3) -- (0,3) node[above] {$y$};
			
			% 椭圆 e=0.6
			\draw[blue, thick] (1,0) ellipse (1.5 and 1.2);
			\node[blue] at (1, -2) {椭圆 ($e=0.6$)};
			
			% 抛物线 e=1
			\draw[red, thick, domain=-0.5:2.5] plot (\x, {1.5*\x*\x - 1.5});
			\node[red] at (4, 2) {抛物线 ($e=1$)};
			
			% 双曲线 e=1.5
			\draw[green!70!black, thick, domain=-1:1] plot ({2.5+1.5*cosh(\x)}, {1.5*sinh(\x)});
			\draw[green!70!black, thick, domain=-1:1] plot ({2.5-1.5*cosh(\x)}, {1.5*sinh(\x)});
			\node[green!70!black] at (4, -2) {双曲线 ($e=1.5$)};
		\end{tikzpicture}
	\end{center}
	
	\subsection*{三、标准方程推导（以焦点在$x$轴为例）}
	
	建立直角坐标系：
	\begin{itemize}
		\item 设焦点 \( F = (p, 0) \)
		\item 设准线 \( l: x = -p \)
		\item 设动点 \( P = (x, y) \)
	\end{itemize}
	
	由统一定义可得：
	\[
	\frac{\sqrt{(x - p)^2 + y^2}}{|x + p|} = e
	\]
	
	两边平方并整理，得到统一方程：
	\[
	(1 - e^2)x^2 + y^2 - 2p(1 + e^2)x + p^2(1 - e^2) = 0
	\]
	
	\subsubsection*{3.1 椭圆 (\( 0 < e < 1 \))}
	令 \( a = \dfrac{p e}{1 - e^2} \), \( c = ae \)，方程可化为标准形式：
	\[
	\frac{(x - h)^2}{a^2} + \frac{y^2}{b^2} = 1 \quad \text{其中 } b^2 = a^2 - c^2
	\]
	此时准线方程为 \( x = \pm \dfrac{a^2}{c} \)。
	
	\subsubsection*{3.2 抛物线 (\( e = 1 \))}
	方程简化为标准形式：
	\[
	y^2 = 4px
	\]
	
	\subsubsection*{3.3 双曲线 (\( e > 1 \))}
	令 \( a = \dfrac{p e}{e^2 - 1} \), \( c = ae \)，方程可化为标准形式：
	\[
	\frac{(x - h)^2}{a^2} - \frac{y^2}{b^2} = 1 \quad \text{其中 } b^2 = c^2 - a^2
	\]
	此时准线方程为 \( x = \pm \dfrac{a^2}{c} \)。
	
	\subsection*{四、教学说明与几何意义}
	
	\begin{enumerate}
		\item \textbf{统一性}：此定义将三类圆锥曲线统一于一个简洁的几何条件，揭示了它们的内在联系。
		
		\item \textbf{历史渊源}：此定义与“用平面截割圆锥面”的生成方式完全等价，是其在平面几何中的完美表述。
		
		\item \textbf{离心率 \( e \)} 的几何意义：
		\begin{itemize}
			\item 是描述曲线形状的\textbf{核心参数}
			\item \( e \) 越大，椭圆越扁，双曲线开口越开阔
			\item \( e \to 0 \) 时，椭圆趋近于圆；\( e \to \infty \) 时，双曲线趋近于两条平行线
		\end{itemize}
	\end{enumerate}


\section{圆锥曲线大题}
	
	圆锥曲线大题（尤其解析几何压轴）确实是高考数学里``分数上限''和``思维深度''最明显的板块。想要从不稳到稳拿高分，必须掌握\textbf{结构化思考方法}和\textbf{几何-代数混合策略}。下面我给你一个硬核分析，分为三个层次：\textbf{本质规律 $\rightarrow$ 答题流程模板 $\rightarrow$ 实战技巧与错误陷阱}。
	
	
	\section{本质规律：圆锥曲线题考的不是公式，而是``几何建模 + 代数结构控制力''}
	
	所有高考圆锥曲线大题，无论是椭圆、双曲线、抛物线，背后都考这三类能力：
	
		\begin{enumerate}[label=\arabic*.]
			\item \textbf{几何建模能力}
			\par —— 把文字条件转化为点、线、动点轨迹的代数约束。
			\par 比如：``点 P 在椭圆上且满足 $PF_1 = 2PF_2$''，立刻写出椭圆方程 + 距离公式建立方程组。
			
			\item \textbf{代数结构控制能力}
			\par —— 用参数法、消元、判别式、取最值等手段``降维''。
			\par 如：设动点坐标 $(x, y)$，转化成一个单变量函数，最后借助判别式或一元二次方程解的存在性解决问题。
			
			\item \textbf{数形结合与极值意识}
			\par —— 关键突破口几乎都来自``图像关系''或``几何意义''。
			\par 比如：``点到定直线距离最小''，你若代入距离公式死算，极容易错；换成几何对称反射思想，立刻变简单。
		\end{enumerate}

	
	\section{答题流程模板：几乎所有圆锥曲线压轴题都能套进这 5 步}
	
		\begin{table}[ht]
		\centering
		\begin{tabular}{|p{2.5cm}|p{4cm}|p{5.5cm}|}
			\hline
			\textbf{步骤} & \textbf{核心任务} & \textbf{常见技巧} \\
			\hline
			\textbf{1. 建系定参} & 先定坐标系，确定参数 $a$、$b$、$e$。 & 若题中提及``焦点''、``准线''、``离心率''等，\textbf{优先选标准方程}。否则，可用一般式 $Ax^2 + Bxy + Cy^2 + Dx + Ey + F = 0$。 \\
			\hline
			\textbf{2. 写出条件方程} & 将动点或线段的几何条件翻译成代数式。 & 常用套路：点在线上 $\rightarrow$ 代入直线方程；到焦点距离关系 $\rightarrow$ 根号公式；共圆/共线 $\rightarrow$ 行列式 = 0。 \\
			\hline
			\textbf{3. 化简/参数化} & 引入参数使问题降维。 & 典型例子：椭圆 $\frac{x^2}{a^2} + \frac{y^2}{b^2} = 1$ 可设 $x = a\cos\theta$, $y = b\sin\theta$；双曲线设 $x = a\sec\theta$, $y = b\tan\theta$。 \\
			\hline
			\textbf{4. 判定与优化} & 研究方程有解条件、最值、范围。 & 主要武器：判别式、均值不等式、导数或几何极值。 \\
			\hline
			\textbf{5. 回代与解释} & 代回原式并说明几何意义。 & 常出错点：没有说明``取等时的几何条件''，导致逻辑不完整。 \\
			\hline
		\end{tabular}
	\end{table}
	

	
	\section{实战技巧与潜规则}
	
	\subsection{动点轨迹型题的万能策略}
	
	常见问法：``点 P 的轨迹是什么？''
	\textbf{万能思路}：
	
	\begin{itemize}
		\item 先把条件化为关于 $x$, $y$ 的方程；
		\item \textbf{若含参数}（如点在直线上），$\rightarrow$ 消去参数得二次方程 $\rightarrow$ 用``判别式恒成立''判断轨迹；
		\item \textbf{若条件含绝对值或距离}，$\rightarrow$ 平方去根 $\rightarrow$ 化成二次曲线标准式。
	\end{itemize}
	
	\textbf{典型例题结构}：
	
	\begin{quote}
		已知椭圆 $\frac{x^2}{a^2}+\frac{y^2}{b^2}=1$，点 $M(x_1,y_1)$ 在椭圆上，点 $P$ 满足 $|PM|+|PF_1|=|PF_2|$，求 P 的轨迹。
		\newline$\rightarrow$ 先写椭圆参数式 $M(a\cos\theta,b\\sin\theta)$，
		\newline$\rightarrow$ 代入距离关系，平方、化简、消参，
		\newline$\rightarrow$ 结果多为``另一条圆锥曲线''。
	\end{quote}
	
	\subsection*{(2) ``动线交点''型题的隐藏突破口}
	
	当题中出现``直线 l 过定点/定斜率/定截距''，并求与圆锥曲线交点条件时：
	
	\begin{itemize}
		\item 写直线方程 $y=kx+b$ 或 $y=k(x-x_0)+y_0$；
		\item 代入圆锥曲线方程 $\rightarrow$ 得到二次方程；
		\item \textbf{利用判别式 $\Delta$}，从``有实交点''或``有重根''条件出发；
		\item 若问切线 $\rightarrow$ 令 $\Delta = 0$；
		\item 若问弦长、面积、P点运动范围 $\rightarrow$ 研究 $\Delta>0$ 的参数区间。
	\end{itemize}
	
	这是最常见的高考套路：
	
	\begin{quote}
		``定点 + 定值 + 切线判别式恒等式''。
	\end{quote}
	
	\subsection*{(3) ``几何+代数双解法''原则}
	
	压轴题如果你\textbf{只会算}，很可能算崩。要始终问自己：
	
	\begin{quote}
		有没有几何图像可以简化关系？
	\end{quote}
	
	例如：
	
	\begin{itemize}
		\item \textbf{最短距离问题}：用``对称法''；
		\item \textbf{轨迹问题}：用``焦点定义或反射性质''；
		\item \textbf{弦中点、垂直关系}：用``极点极线''思想（即利用几何对偶性，减少代数操作）。
	\end{itemize}
	
	\subsection*{(4) 常用结论速记（建议熟背）}
	
		\begin{table}[ht]
		\centering
		\begin{tabular}{|p{3cm}|p{8cm}|}
			\hline
			\textbf{类型} & \textbf{结论} \\
			\hline
			椭圆焦点反射性质 & 光从焦点射出，经椭圆反射会经过另一焦点。$\rightarrow$ 常用于``角平分线/最值问题''。 \\
			\hline
			椭圆上一点 P 切线方程 & $\frac{x_1x}{a^2} + \frac{y_1y}{b^2} = 1$。 \\
			\hline
			双曲线切线方程 & $\frac{x_1x}{a^2} - \frac{y_1y}{b^2} = 1$。 \\
			\hline
			抛物线切线方程 & $y_1y = 2p(x+x_1)$。 \\
			\hline
			椭圆弦中点坐标 & 若弦两端点参数为 $\theta_1$, $\theta_2$，则中点参数为 $\frac{\theta_1+\theta_2}{2}$。 \\
			\hline
			椭圆到直线距离最小值 & 由法线方程 + 几何反射法求，非直接代数最值。 \\
			\hline
		\end{tabular}
	\end{table}
	
	\subsection*{(5) 高频错误点（要强制规避）}
	
	\begin{enumerate}[label=\arabic*.]
		\item \textbf{忽略定义域/范围条件}（如令判别式$\geq 0$后忘记限制参数范围）。
		\item \textbf{解得轨迹后没检查合法性}（例如平方去根造成虚解）。
		\item \textbf{表达不严谨}（尤其是``若…则…''、``当且仅当…''的逻辑关系不写）。
		\item \textbf{画图不准}：高分同学都会画标准椭圆辅助图，标焦点、准线、动点方向、切线、斜率符号。
	\end{enumerate}
	

	
	\section{圆锥曲线压轴题的``心法''}
	
	\begin{center}
		\fbox{\parbox{0.9\textwidth}{
				\centering
				\textbf{本质：几何条件代数化 $\rightarrow$ 参数消元降维 $\rightarrow$ 判别式或最值处理 $\rightarrow$ 几何解释。}
		}}
	\end{center}
	
	如果你能把所有题都拆解为``一个几何约束 + 一个代数条件''，那么圆锥曲线大题就再也不是黑箱。
	
	
	\begin{table}[ht]
		\centering
		\begin{tabular}{|c|c|c|c|}
			\hline
			\textbf{椭圆} & 焦点三角形 & 面积公式：$S_{\triangle PF_1F_2} = b^2 \tan \frac{\theta}{2}$ & \\
			\hline
			& 周角定理 & $k_{PA} \cdot k_{PB} = e^2 - 1$ (A, B为长轴端点) & \\
			\hline
			& 中点弦斜率 & $k_{AB} \cdot k_{OM} = e^2 - 1$ (M为弦AB中点) & \\
			\hline
			& 焦点弦与离心率 & 已知倾斜角$\alpha$且$\|\overrightarrow{AF}\| = \lambda \|\overrightarrow{FB}\|$，则$e=\left|\frac{\lambda-1}{(\lambda+1)\cos\alpha}\right|$ & \\
			\hline
		\end{tabular}
	\end{table}
	
	\begin{table}[ht]
		\centering
		\begin{tabular}{|c|c|c|c|}
			\hline
			\textbf{双曲线} & 焦点三角形 & 面积公式：$S_{\triangle PF_1F_2} = \frac{b^2}{\tan(\theta/2)}$ & \\
			\hline
			& 周角定理 & $k_{PA} \cdot k_{PB} = e^2 - 1$ (A, B为实轴端点) & \\
			\hline
			& 中点弦斜率 & $k_{AB} = \frac{b^2 x_0}{a^2 y_0}$ (M$(x_0,y_0)$为弦AB中点) & \\
			\hline
			& 焦点弦与离心率 & 交于右支且$\|\overrightarrow{AF}\| = \lambda \|\overrightarrow{FB}\|$，则$e=\left|\frac{\lambda-1}{(\lambda+1)\cos\alpha}\right|$ & \\
			\hline
		\end{tabular}
	\end{table}
	
		\begin{table}[ht]
			\centering
			\begin{tabular}{|p{1.5cm}|p{2.4cm}|p{8.8cm}|}
				\hline
				\textbf{抛物线} & 焦点弦长 & $\|AB\| = \frac{2p}{\sin^2 \theta}$ ($\theta$为倾斜角) \\
				\hline
				& 焦半径倒数和 & $\frac{1}{\|AF\|} + \frac{1}{\|BF\|} = \frac{2}{p}$ \\
				\hline
				& 焦点弦性质 & 以AB为直径的圆与准线相切；A, O, D三点共线（D为准线与轴交点） \\
				\hline
			\end{tabular}
		\end{table}
		
		\begin{table}[ht]
			\centering
			\begin{tabular}{|p{1.5cm}|p{2.4cm}|p{8.8cm}|}
				\hline
				\textbf{综合性} & 切点弦方程 & 曲线外一点P$(x_0,y_0)$引两条切线，切点弦AB所在直线方程（椭圆为例）：$\frac{x_0 x}{a^2} + \frac{y_0 y}{b^2} = 1$ \\
				\hline
				& 定点定直线 & 圆锥曲线内接直角三角形的直角顶点与圆锥曲线的顶点重合，则斜边所在直线过定点 \\
				\hline
				& 斜率积定值 & 过定点作两直线与曲线交于两点，若两直线斜率之积为定值，则两交点连线恒过定点 \\
				\hline
			\end{tabular}
	\end{table}
	
	\section*{结论详解与运用提示}
	
	\begin{itemize}
		\item \textbf{焦点三角形面积公式（椭圆）}：在椭圆中，若$\angle F_1PF_2 = \theta$，则面积$S_{\triangle PF_1F_2} = b^2 \tan \frac{\theta}{2}$。双曲线中为$S_{\triangle PF_1F_2} = \frac{b^2}{\tan(\theta/2)}$，注意区分。
		
		\item \textbf{周角定理}：在椭圆（或双曲线）中，若点P在曲线$C$上，A, B为椭圆长轴（或双曲线实轴）端点，则有$k_{PA} \cdot k_{PB} = e^2 - 1$。在椭圆中，$e^2 - 1 = -\frac{b^2}{a^2}$，故斜率积为负值；在双曲线中，$e^2 - 1 = \frac{b^2}{a^2}$，故斜率积为正值。
		
		\item \textbf{中点弦斜率关系（有心圆锥曲线垂径定理）}：适用于椭圆和双曲线。对于标准方程$C: \frac{x^2}{a^2} \pm \frac{y^2}{b^2} = 1$，弦AB中点M$(x_0, y_0)$，则$k_{AB} \cdot k_{OM} = \mp \frac{b^2}{a^2} = e^2 - 1$。圆也符合此规律。
		
		\item \textbf{抛物线焦点弦}：设倾斜角为$\theta$，焦点弦长$|AB| = \frac{2p}{\\sin^2 \theta}$。另外，$\frac{1}{|AF|} + \frac{1}{|BF|} = \frac{2}{p}$也是一个常用结论。
		
		\item \textbf{切点弦方程}：过圆锥曲线外一点P$(x_0, y_0)$作两条切线，切点分别为A, B，则切点弦AB所在直线方程可通过将曲线方程中的$x^2$替换为$x_0x$，$y^2$替换为$y_0y$等方式快速得到。例如，椭圆$\frac{x^2}{a^2} + \frac{y^2}{b^2} = 1$的切点弦方程为$\frac{x_0 x}{a^2} + \frac{y_0 y}{b^2} = 1$。
	\end{itemize}
	

   圆锥曲线的统一方程在极坐标和直角坐标系中都有表述，以下是常用的形式
   
    极坐标下的统一方程
   
   最常用的统一方程为
   
  
   $ r = \frac{ep}{1 + e\cos\theta}$
   
   
   或者等价的
   
   $r = \frac{p}{1 + e\cos\theta} $
   
   其中：
   - $r$ 为极径
   - $\theta$ 为极角
   - $e$ 为离心率
   - $p$ 为焦参数（焦点到准线的距离）
   
    直角坐标系中的统一形式
   
   在直角坐标系中，圆锥曲线可以统一表示为：
   
   
   $Ax^2 + Bxy + Cy^2 + Dx + Ey + F = 0$
   
   
   其中 $A, B, C$ 不同时为0。
   
    根据离心率的分类
   
%   
%   \begin{cases}
%   	\text{椭圆：} $& 0 \leq e < 1 $
%   	\text{抛物线：}$ & e = 1 $
%   	\text{双曲线：} $ & e > 1 $
%   \end{cases}
%   
   
   标准位置的统一方程
   
   当焦点在原点，准线为 $x = p$ 时：
   
  
  $ \sqrt{x^2 + y^2} = e(p - x)$
   
   
    重要说明
   
   - 极坐标形式最为简洁统一，能够涵盖所有类型的圆锥曲线
   - 直角坐标形式是隐式方程，需要通过判别式 $B^2 - 4AC$ 来判断曲线类型
   - 这些统一方程在解决涉及离心率的问题时特别有用
   

   
   \section{阿波罗尼斯圆}
   
   \textbf{学习目标}
   \begin{itemize}
   	\item 理解阿波罗尼斯圆的\textbf{定义}与\textbf{几何直观}。
   	\item 掌握\textbf{方程推导}、\textbf{圆心位置与半径公式}、\textbf{退化情形}。
   	\item 能用阿波罗尼斯圆处理\textbf{比例轨迹}与\textbf{三角形定比问题}。
   \end{itemize}
   
   \textbf{一、定义与直观}
   \begin{itemize}
   	\item 给定平面两定点 \(A,B\) 与正数 \(k>0\)。满足
   	\[
   	\frac{PA}{PB}=k \qquad (P \text{ 为动点})
   	\]
   	的点的轨迹是一条\textbf{圆}，称为以 \(A,B\) 为焦点、比值为 \(k\) 的\textbf{阿波罗尼斯圆}。当 \(k=1\) 时，轨迹为线段 \(\overline{AB}\) 的\textbf{垂直平分线}（退化为直线）。
   	\item 与“定和为常”的椭圆类比：此处为\(\textbf{定比为常}\)，因而得到的是圆。
   \end{itemize}
   
   \textbf{二、代数推导（圆心与半径）}
   \begin{itemize}
   	\item 取坐标系使 \(A(0,0),\,B(d,0)\)（\(d=AB>0\)），令 \(P(x,y)\) 满足 \(\lvert PA\rvert=k\lvert PB\rvert\)（\(k\neq1\)）。则
   	\begin{align*}
   		x^2+y^2 &= k^2\bigl((x-d)^2+y^2\bigr)\\
   		\Longleftrightarrow\quad (x-\tfrac{k^2}{k^2-1}d)^2+y^2 &= \left(\tfrac{k}{\lvert k^2-1\rvert}d\right)^2.
   	\end{align*}
   	因此阿波罗尼斯圆的\textbf{标准形式}为
   	\[
   	\boxed{\ (x-\tfrac{k^2}{k^2-1}d)^2+y^2=\left(\tfrac{k}{\lvert k^2-1\rvert}d\right)^2\ }.
   	\]
   	\item \textbf{圆心}位于直线 \(AB\) 上：
   	\[
   	O\Bigl(\tfrac{k^2}{k^2-1}d,\,0\Bigr),\qquad
   	\frac{OA}{OB}=k^2\ \text{（定向分比，可能为外分）}.
   	\]
   	\item \textbf{半径}
   	\[
   	r=\frac{k}{\lvert k^2-1\rvert}\,AB.
   	\]
   	\item \textbf{与 \(AB\) 的交点}：令 \(y=0\) 得
   	\[
   	x=\tfrac{k^2}{k^2-1}d\ \pm\ \tfrac{k}{\lvert k^2-1\rvert}d,
   	\]
   	两点分别为 \(\overline{AB}\) 的\textbf{内分点}与\textbf{外分点}（分比 \(k:1\)）。
   	\item \textbf{退化情形}：\(k=1\Rightarrow\) 轨迹为 \(\overline{AB}\) 的垂直平分线；\(k\to0^+\) 或 \(k\to+\infty\) 时，圆半径趋于 \(0\)，轨迹分别向点 \(A\) 或 \(B\) 收缩。
   \end{itemize}
   
   \textbf{三、经典二级结论（必备）}
   \begin{enumerate}
   	\item \textbf{定点定比定圆}：给定 \(A,B,k(\neq1)\) 唯一确定一条阿波罗尼斯圆。
   	\item \textbf{分点性质}：阿波罗尼斯圆与 \(AB\) 的交点是 \(\overline{AB}\) 的内分/外分点，分比为 \(k:1\)。
   	\item \textbf{到定点距离的范围}：设 \(d=AB\)，则
   	\[
   	OA=\frac{k^2}{\lvert k^2-1\rvert}d,\qquad r=\frac{k}{\lvert k^2-1\rvert}d.
   	\]
   	故 \(\lvert PA\rvert\in\bigl[\lvert OA-r\rvert,\ OA+r\bigr]\)，且 \(\lvert PB\rvert=\lvert PA\rvert/k\)。
   	\item \textbf{三角形中的 A-阿波罗尼斯圆}：在 \(\triangle ABC\) 中，满足
   	\[
   	\frac{PB}{PC}=\frac{AB}{AC}=:k
   	\]
   	的点 \(P\) 构成一圆（以 \(B,C\) 为焦点，比例 \(k\)）。其\textbf{圆心}在直线 \(BC\) 上，且
   	\[
   	\frac{OB}{OC}=k^2=\Bigl(\frac{AB}{AC}\Bigr)^2,\qquad
   	r=\frac{k}{\lvert k^2-1\rvert}\,BC,
   	\]
   	并且该圆与 \(BC\) 的交点正是按 \(AB:AC\) 对 \(BC\) 的\textbf{内分点}与\textbf{外分点}；点 \(A\) 在此圆上（因 \( \tfrac{AB}{AC}=k \)）。
   	\item \textbf{同轴性}：对固定的 \(A,B\)，不同 \(k\) 的阿波罗尼斯圆与 \(\overline{AB}\) \textbf{同轴}（圆心皆在直线 \(AB\) 上）。
   \end{enumerate}
   
   \textbf{四、典型例题（含解析）}
   
   \textbf{例1（方程与要素）}\quad 已知 \(A(0,0),\,B(4,0)\)。求轨迹 \(\dfrac{PA}{PB}=2\) 的方程、圆心、半径及与 \(AB\) 的交点。
   \begin{itemize}
   	\item 代入通式（\(d=4,\,k=2\)）：
   	\[
   	(x-\tfrac{2^2}{2^2-1}\cdot4)^2+y^2=\left(\tfrac{2}{2^2-1}\cdot4\right)^2
   	\ \Longrightarrow\
   	(x-\tfrac{16}{3})^2+y^2=\left(\tfrac{8}{3}\right)^2.
   	\]
   	\item 圆心 \(O\!\left(\tfrac{16}{3},\,0\right)\)，半径 \(r=\tfrac{8}{3}\)。
   	\item 与 \(AB\) 的交点（令 \(y=0\)）：
   	\[
   	x=\tfrac{16}{3}\pm\tfrac{8}{3}\Rightarrow x=\tfrac{8}{3},\ 8.
   	\]
   	其中 \(x=\tfrac{8}{3}\) 为\textbf{内分点}（分比 \(2:1\)），\(x=8\) 为\textbf{外分点}。
   \end{itemize}
   
   \textbf{例2（到定点距离的范围）}\quad 设 \(AB=10\)，动点 \(P\) 满足 \(\dfrac{PA}{PB}=2\)。求 \(\lvert PA\rvert,\lvert PB\rvert\) 的取值范围及最值。
   \begin{itemize}
   	\item \(OA=\dfrac{2^2}{2^2-1}\cdot10=\dfrac{40}{3}\),\quad \(r=\dfrac{2}{2^2-1}\cdot10=\dfrac{20}{3}\)。
   	\item \(\lvert PA\rvert\in\bigl[OA-r,\,OA+r\bigr]=\bigl[\tfrac{20}{3},\,20\bigr]\)；\quad
   	\(\lvert PB\rvert=\lvert PA\rvert/2\in\bigl[\tfrac{10}{3},\,10\bigr]\)。
   	\item 最小值在靠近 \(A\) 的交点处取得，最大值在远离 \(A\) 的交点处取得（均在直线 \(AB\) 上）。
   \end{itemize}
   
   \textbf{例3（三角形中的 A-阿波罗尼斯圆）}\quad 在 \(\triangle ABC\) 中，令 \(k=\dfrac{AB}{AC}\)。证明：满足 \(\dfrac{PB}{PC}=k\) 的点 \(P\) 构成一圆，其圆心在 \(BC\) 上，并求该圆与 \(BC\) 的交点。
   \begin{itemize}
   	\item 证：以 \(B,C\) 为焦点、比值 \(k\) 的阿波罗尼斯圆之定义即给出轨迹为圆；由二级结论，
   	\(\dfrac{OB}{OC}=k^2\)，圆心在 \(BC\) 上。
   	\item 与 \(BC\) 的交点为按 \(k:1=\dfrac{AB}{AC}:1\) 对 \(BC\) 的\textbf{内分点}与\textbf{外分点}；且 \(A\) 在该圆上（因 \(\dfrac{AB}{AC}=k\)）。
   \end{itemize}
   
   \textbf{五、同步练习（自测）}
   \begin{enumerate}
   	\item 已知 \(A(-2,0),\,B(3,0)\)。求轨迹 \(\dfrac{PA}{PB}=\dfrac{3}{2}\) 的方程、圆心与半径，并写出与 \(AB\) 的交点坐标。
   	\item 在 \(\triangle ABC\) 中，设 \(AB:AC=5:3\)。写出 A-阿波罗尼斯圆的 \(\dfrac{OB}{OC}\) 与半径关于 \(BC\) 的表达式，并判断点 \(A\) 是否在该圆上（说明理由）。
   \end{enumerate}
   
   \textbf{要点小结}
   \begin{itemize}
   	\item \(\dfrac{PA}{PB}=k\)（\(k>0\)）的轨迹为圆；\(k=1\) 退化为垂直平分线。
   	\item 圆心在 \(AB\) 上，满足 \(\dfrac{OA}{OB}=k^2\)；半径 \(r=\dfrac{k}{\lvert k^2-1\rvert}AB\)。
   	\item 与 \(AB\) 的两交点为分比 \(k:1\) 的内分/外分点；三角形中满足 \(\dfrac{PB}{PC}=\dfrac{AB}{AC}\) 的点集为 A-阿波罗尼斯圆（心在 \(BC\) 上，过点 \(A\)）。
   \end{itemize}
   
   \section{非对称问题}
   
   已知椭圆 \(\dfrac{x^{2}}{4}+\dfrac{y^{2}}{3}=1\) 的左右顶点分别为 \(A,B\)，过椭圆的右焦点 \(F(1,0)\) 的直线 \(l\) 与椭圆交于 \(C,D\) 两点（不与 \(A,B\) 重合）。记直线 \(AC\) 与 \(BD\) 的斜率分别为 \(k_{1},k_{2}\)，证明：\(\dfrac{k_{1}}{k_{2}}\) 为定值
   
   \textbf{Corrected solution.}
   椭圆
   \[
   \frac{x^{2}}{4}+\frac{y^{2}}{3}=1
   \]
   的左顶点为 \(A(-2,0)\)，右顶点为 \(B(2,0)\)，右焦点为 \(F(1,0)\)。
   
   设过 \(F(1,0)\) 的直线斜率为 \(m\)，其方程为
   \[
   y = m(x-1).
   \]
   把它代入椭圆方程：
   \[
   \frac{x^{2}}{4} + \frac{m^{2}(x-1)^{2}}{3} = 1,
   \]
   两边乘以 \(12\) 得
   \[
   3x^{2} + 4m^{2}(x-1)^{2} = 12,
   \]
   即
   \[
   (3+4m^{2})x^{2} - 8m^{2}x + (4m^{2}-12) = 0. \tag{1}
   \]
   这是直线与椭圆的两个交点 \(C,D\) 的 \(x\) 坐标方程。
   
   判别式
   \[
   \Delta = (-8m^{2})^{2} - 4(3+4m^{2})(4m^{2}-12)
   = 144(m^{2}+1),
   \]
   所以
   \[
   x_{C,D}
   = \frac{8m^{2} \pm 12\sqrt{m^{2}+1}}{2(3+4m^{2})}
   = \frac{4m^{2} \pm 6\sqrt{m^{2}+1}}{4m^{2}+3}.
   \]
   记
   \[
   u = \sqrt{m^{2}+1}, \quad
   x_C = \frac{4m^{2}+6u}{4m^{2}+3}, \quad
   x_D = \frac{4m^{2}-6u}{4m^{2}+3},
   \]
   并有
   \[
   C\bigl(x_C,\, m(x_C-1)\bigr), \qquad
   D\bigl(x_D,\, m(x_D-1)\bigr).
   \]
   
   因为本题约定的是“左顶点为 \(A\)，右顶点为 \(B\)”，所以
   \[
   A(-2,0), \quad B(2,0).
   \]
   
   于是
   \[
   k_1 = \text{slope of } AC
   = \frac{m(x_C-1)-0}{x_C - (-2)}
   = m \cdot \frac{x_C - 1}{x_C + 2},
   \]
   \[
   k_2 = \text{slope of } BD
   = \frac{m(x_D-1)-0}{x_D - 2}
   = m \cdot \frac{x_D - 1}{x_D - 2}.
   \]
   所以
   \[
   \frac{k_1}{k_2}
   = \frac{m \dfrac{x_C - 1}{x_C + 2}}{m \dfrac{x_D - 1}{x_D - 2}}
   = \frac{(x_C - 1)(x_D - 2)}{(x_C + 2)(x_D - 1)}. \tag{2}
   \]
   
   下面把 \(x_C, x_D\) 代入。先写出四个差：
   \[
   x_C - 1 = \frac{6u - 3}{4m^{2}+3}, \quad
   x_C + 2 = \frac{12m^{2} + 6u + 6}{4m^{2}+3},
   \]
   \[
   x_D - 1 = \frac{-6u - 3}{4m^{2}+3}, \quad
   x_D - 2 = \frac{-4m^{2} - 6u - 6}{4m^{2}+3}.
   \]
   代入 (2) 得
   \[
   \frac{k_1}{k_2}
   = \frac{\dfrac{6u - 3}{4m^{2}+3} \cdot \dfrac{-4m^{2} - 6u - 6}{4m^{2}+3}}
   {\dfrac{12m^{2} + 6u + 6}{4m^{2}+3} \cdot \dfrac{-6u - 3}{4m^{2}+3}}
   = \frac{(6u - 3)(-4m^{2} - 6u - 6)}{(12m^{2} + 6u + 6)(-6u - 3)}.
   \]
   把公因子提出来并约掉，可以化成
   \[
   \frac{k_1}{k_2} = \frac{1}{3} \cdot
   \frac{(2u - 1)(2m^{2} + 3u + 3)}{(2m^{2} + u + 1)(2u + 1)}.
   \]
   注意到 \(u^{2} = m^{2} + 1\)，直接展开分子、分母，可以验证
   \[
   (2u - 1)(2m^{2} + 3u + 3) = (2m^{2} + u + 1)(2u + 1),
   \]
   因此上式中这部分恰好等于 1，故
   \[
   \boxed{\frac{k_1}{k_2} = \frac{1}{3}}
   \]
   与斜率 \(m\) 无关，是一个定值。
   
   
   \medskip
   
   \textbf{练习}：已知椭圆
   \[
   C:\ \frac{x^{2}}{a^{2}}+\frac{y^{2}}{b^{2}}=1 \quad (a>b>0)
   \]
   的左、右焦点分别为 \(F_{1}(-c,0)\)、\(F_{2}(c,0)\)，\(M,N\) 分别为左、右顶点，直线
   \[
   l:\ x=ty+1
   \]
   与椭圆 \(C\) 交于 \(A,B\) 两点，当 \(t=-\dfrac{\sqrt{3}}{3}\) 时，\(A\) 是椭圆的上顶点，且 \(\triangle AF_{1}F_{2}\) 的周长为 \(6\)。
   
   (1) 求椭圆 \(C\) 的方程；
   
   (2) 设直线 \(AM, BN\) 交于点 \(Q\)，证明：点 \(Q\) 在定直线上；
   
   (3) 设直线 \(AM, BN\) 的斜率分别为 \(k_{1},k_{2}\)，证明：\(\dfrac{k_{1}}{k_{2}}\) 为定值。
   
   
   \section{记录}
   
   \subsection{Holder不等式}
   
   \subsection{Carlson不等式}
   
 
   Cauchy不等式
   \[
   (a_1 + b_1)(a_2 + b_2) \geq  ( \sqrt{a_1 a_2} + \sqrt{b_1 b_2})^2 
   \]

	Karlson不等式
	
	\[
	(a_1 + b_1)(a_2 + b_2)(a_3 + b_3) \geq  ( \sqrt[3]{a_1 a_2 a_3} + \sqrt[3]{b_1 b_2 b_3})^3
	\]
	
	\begin{example}
		  如果\( \theta \in (0, \frac{\pi}{2} ), \) 则\( \frac{8}{\sin \theta} + \frac{1}{cos \theta} \) 的最小值为 \( 5 \sqrt{5} \)
		
	\end{example}
	
	\begin{proof}
		\begin{equation*}
			(\frac{8}{\sin \theta} + \frac{1}{cos \theta})( \frac{8}{\sin \theta} + \frac{1}{cos \theta}) (\sin^2 \theta + \cos^2 \theta) \geq ()^3
		\end{equation*}
	\end{proof}
	
	
	\section{射影几何与极点极线理论}

\section{圆锥曲线六大定义}

\subsection{}

\section{焦半径公式与焦点弦定理}

\section{椭圆双曲线焦点三角形命题点}

\section{三一线模型}

\section{抛物线25条结论推导}

\section{圆锥曲线光学性质与工程技术应用}

\section{点乘双根}

\textbf{适用类型：} 类似 \( x_1, x_2, y_1, y_2 \)，\( (x_1 + t)(x_2 + t) \)，\( (y_1 + t)(y_2 + t) \) 或 \( MA \cdot MB \)，\( |MA| \cdot |MB| \)（其中 \( x_1, x_2, y_1, y_2 \) 是直线与曲线的两个交点的横纵坐标，\( A, B \) 直线与曲线的两个交点）以及可转化为上述结构的问题。

\textbf{理论基础：} 二次函数的双根式，若一元二次方程 \( ax^2 + bx + c = 0 \) (\( a \neq 0 \)) 的两根为 \( x_1, x_2 \)，则
\[
ax^2 + bx + c = a(x - x_1)(x - x_2)
\]

\textbf{具体步骤：} 化双根式 $\rightarrow$ 赋值 $\rightarrow$ 变形代入

\noindent 1. 椭圆 \( C: \frac{x^2}{4} + \frac{y^2}{3} = 1 \)，左顶点为 \( A \)，过左焦点 \( P \) 的直线交椭圆于 \( P, Q \) 两点，求证：直线 \( AP \) 和 \( AQ \) 的斜率乘积是定值。
	

\textbf{题目：}椭圆
\[
C:\ \frac{x^{2}}{4}+\frac{y^{2}}{3}=1,
\]
左顶点为 $A$，过左焦点 $F$ 的直线交椭圆于 $P,Q$ 两点，求证：直线 $AP, AQ$ 的斜率乘积为定值。

\textbf{解：}由椭圆方程可得
\[
a=2,\quad b=\sqrt3,\quad c=\sqrt{a^{2}-b^{2}}=1,
\]
故左顶点 $A(-2,0)$，左焦点 $F(-1,0)$。

设过焦点 $F$ 的一条动直线斜率为 $t\ (\neq0)$，则其方程为
\[
y=t(x+1). \tag{1}
\]
该直线与椭圆相交于 $P(x_1,y_1),Q(x_2,y_2)$ 两点，将 (1) 代入椭圆方程
\[
\frac{x^{2}}{4}+\frac{t^{2}(x+1)^{2}}{3}=1
\]
化简得
\[
\frac{1}{12}\Big[(4t^{2}+3)x^{2}+8t^{2}x+(4t^{2}-12)\Big]=0,
\]
根不受常数因子影响，于是 $x_1,x_2$ 为二次方程
\[
(4t^{2}+3)x^{2}+8t^{2}x+(4t^{2}-12)=0 \tag{2}
\]
的两根。

对任一点 $M(x,y)$（为 $P$ 或 $Q$），由 (1) 得 $y=t(x+1)$，直线 $AM$ 的斜率为
\[
k=\frac{y-0}{x-(-2)}=\frac{t(x+1)}{x+2}.
\]
因此
\[
k_{AP}=\frac{t(x_1+1)}{x_1+2},\qquad
k_{AQ}=\frac{t(x_2+1)}{x_2+2},
\]
从而
\[
k_{AP}\cdot k_{AQ}
=t^{2}\frac{(x_1+1)(x_2+1)}{(x_1+2)(x_2+2)}. \tag{3}
\]

下面用“点乘双根法”计算分子、分母。

由 (2) 可写成
\[
f(x)=(4t^{2}+3)x^{2}+8t^{2}x+(4t^{2}-12)
=(4t^{2}+3)(x-x_1)(x-x_2).
\]
令 $x=-1$，得到
\[
f(-1)=(4t^{2}+3)(-1-x_1)(-1-x_2)=(4t^{2}+3)(x_1+1)(x_2+1),
\]
即
\[
(x_1+1)(x_2+1)=\frac{f(-1)}{4t^{2}+3}.
\]
同理令 $x=-2$，
\[
f(-2)=(4t^{2}+3)(-2-x_1)(-2-x_2)=(4t^{2}+3)(x_1+2)(x_2+2),
\]
故
\[
(x_1+2)(x_2+2)=\frac{f(-2)}{4t^{2}+3}.
\]
这正是“点乘双根法”的应用：通过在二次式两端代入同一点，得到关于两根的乘积型表达式。

直接计算：
\[
\begin{aligned}
	f(-1)
	&=(4t^{2}+3)\cdot1+8t^{2}\cdot(-1)+(4t^{2}-12)  \\
	&=4t^{2}+3-8t^{2}+4t^{2}-12
	=-9, \\[4pt]
	f(-2)
	&=(4t^{2}+3)\cdot4+8t^{2}\cdot(-2)+(4t^{2}-12) \\
	&=16t^{2}+12-16t^{2}+4t^{2}-12
	=4t^{2}.
\end{aligned}
\]
于是
\[
(x_1+1)(x_2+1)=\frac{-9}{4t^{2}+3},\qquad
(x_1+2)(x_2+2)=\frac{4t^{2}}{4t^{2}+3}.
\]

代回 (3) 得
\[
k_{AP}\cdot k_{AQ}
=t^{2}\frac{\dfrac{-9}{4t^{2}+3}}{\dfrac{4t^{2}}{4t^{2}+3}}
=t^{2}\cdot\frac{-9}{4t^{2}}
=-\frac{9}{4}.
\]

可见直线 $AP$ 与 $AQ$ 的斜率乘积为
\[
k_{AP}\cdot k_{AQ}=-\frac{9}{4},
\]
与动直线的斜率 $t$ 无关，为一常数。证明完毕。

	\subsection{平移齐次化法}
	
	
		\section{非对称韦达定理}
	

	\section*{圆锥曲线——点乘双根法}

\subsection*{适用类型}

类似
\[
x_1x_2,\qquad y_1y_2,\qquad (x_1+t)(x_2+t),\qquad (y_1+t)(y_2+t),
\]
或
\[
\overrightarrow{MA}\cdot\overrightarrow{MB},\qquad
|\overrightarrow{MA}|\cdot|\overrightarrow{MB}|
\]
的结构。

其中 $x_1,x_2,y_1,y_2$ 是直线与曲线的两个交点的横纵坐标，
$A,B$ 是直线与曲线的两个交点。该方法也适用于可以转化为上述结构的问题。

\subsection*{理论基础}

二次函数的双根式。若一元二次方程
\[
ax^2+bx+c=0\qquad (a\ne 0)
\]
的两根为 $x_1,x_2$，则
\[
ax^2+bx+c=a(x-x_1)(x-x_2).
\]

\subsection*{具体步骤}

\[
\text{化双根式}\longrightarrow \text{赋值}\longrightarrow \text{变形代入}.
\]

\section*{例题 1}

椭圆
\[
C:\frac{x^2}{4}+\frac{y^2}{3}=1
\]
左顶点为 $A$，过左焦点 $F$ 的直线交椭圆于 $P,Q$ 两点。

求证：直线 $AP$ 和 $AQ$ 的斜率乘积是定值。

\section*{例题 2}

（2021 新高考一卷 21 题）

在平面直角坐标系 $xOy$ 中，已知点
\[
F_1(-\sqrt{17},0),\qquad F_2(\sqrt{17},0),
\]
点 $M$ 满足
\[
|MF_1|-|MF_2|=2.
\]
记 $M$ 的轨迹为 $C$。

\begin{enumerate}
  \item 求 $C$ 的方程；
  \[
  x^2-\frac{y^2}{16}=1\qquad (x>0).
  \]

  \item 设点 $T$ 在直线
  \[
  x=\frac12
  \]
  上，过 $T$ 的两条直线分别交 $C$ 于 $A,B$ 两点和 $P,Q$ 两点，且
  \[
  |TA|\cdot|TB|=|TP|\cdot|TQ|.
  \]
  求直线 $AB$ 的斜率与直线 $PQ$ 的斜率之和。
\end{enumerate}

	
	\section{阿基米德三角形}
	\href{https://en.wikipedia.org/wiki/Quadrature_of_the_Parabola}{抛物线的正交性}
\end{document}
